\subsection{The remaining odd triple and two short bare-path eighth moments}\label{proof:Q5-C097}

\paragraph{Q97.0 The triples and their directions.}
Work on a nonempty finite probability space with a real skew kernel
\(|K|\le\allowbreak 1\). The distinct residual terminal arms are \((2,\allowbreak 3,\allowbreak c)\),
\(c\ge\allowbreak 5\); genuine equal pairs may occur at the terminal and receiving
root. The root has an outward unit arm. Every connector edge points
from terminal to root. Choose the terminal direction from its integer
connector length:
\fitdisplay{
\begin{array}{c|c|c}
L&F&\text{ordinary eighth bound}\\\hline
1&\bar f_2h_3f_c&e^{-d/27}\quad\text{(Appendix~\ref{proof:Q5-C070})}\\
\ge2&\bar f_2f_3h_c&e^{-d}
\end{array}
}{}
For \(H=\allowbreak E_rpT^L(E_tF)\), with genuine equal-pair factors \(E_t,\allowbreak E_r\),
we prove the second row:
\begin{equation}\tag{Q97.0.1}\label{eq:Q97:0.1}
\boxed{\E[\{E_rpT^L(E_t\bar f_2f_3h_c)\}^8]\le e^{-d}
\qquad(c\ge5,\ L\ge2).}
\end{equation}
Here \(f_n\) is fully outward, \(\bar f_n\) fully inward, and
\(h_n\) has only its first edge inward. The two directions are selected
by \(L\), before the kernel is given; the low and high domains use the
same nonunit direction. With the first row they yield
\(\E H^8\le\allowbreak  e^{-d/27}\le\allowbreak 1-d/28\) for all \(L\ge\allowbreak 1\).
We also prove, on the whole parameter domain,
\begin{equation}\tag{Q97.0.2}\label{eq:Q97:0.2}
\E[(p\bar f_4)^8]\le e^{3d},\qquad
\boxed{\E[(p\bar f_5)^8]\le e^{-425d/864}.}
\end{equation}

\paragraph{Q97.1 Finite inputs and boundary cases.}
Use the finite definitions of Appendix~\ref{proof:Q5-C064}:
\(A=\allowbreak J+K\), \(T=\allowbreak J-K=\allowbreak A^*\), \(g=\allowbreak K\one\), \(p=\allowbreak 1+g\), \(q=\allowbreak 1-g\),
\(f_n=\allowbreak A^n\one\), \(\bar f_n=\allowbreak T^n\one\), \(h_n=\allowbreak Tf_{n-1}\),
\(\delta_j=\allowbreak \|K^j\one\|_2^2\), \(d=\allowbreak \delta_1\),
\(\theta=\allowbreak 13/32\), \(\kappa=\allowbreak 19/32\), and \(\tau=\allowbreak d-\delta_2\).
All expectations, norms, rows and adjoints use the fixed measure. We have
\begin{equation}\tag{Q97.1.1}\label{eq:Q97:1.1}
\begin{gathered}
0\le d\le1/3,\quad d^2\le\delta_2\le\theta d,\quad
\delta_{j+1}\le\theta\delta_j,\quad\|K\|^2\le\theta,\quad
|Kv(x)|\le\|v\|_2,\\
\E K^{2j}\one=(-1)^j\delta_j,\quad\E K^{2j+1}\one=0,\quad
0\le f_n,\bar f_n,h_n\le2,\quad\E f_n=\E\bar f_n\le1.
\end{gathered}
\end{equation}
The same spectral measure gives \(\delta_3\ge\allowbreak \delta_2^2/d\) for \(d>\allowbreak 0\).
For a mean-zero vector \(v\), the variance of the row
\(K(x,\allowbreak \cdot)-g(x)\) is at most \(1-g(x)^2\), hence
\begin{equation}\tag{Q97.1.2}\label{eq:Q97:1.2}
|Kv(x)|\le\sqrt{p(x)q(x)}\|v\|_2.
\end{equation}
Set \(\beta_1=\allowbreak \beta_2=\allowbreak 0\) and
\(\beta_n=\allowbreak \sum_{j=\allowbreak 1}^{\lfloor(n-1)/2\rfloor}(-1)^{j-1}\delta_j\)
for \(n\ge\allowbreak 3\). Then \(\tau\le\allowbreak \beta_n\le\allowbreak  d\),
\(\beta_3=\allowbreak \beta_4=\allowbreak d\), \(\beta_5=\allowbreak \beta_6=\allowbreak \tau\), and
\begin{equation}\tag{Q97.1.3}\label{eq:Q97:1.3}
S_n=\sum_{j=0}^nK^j\one=f_n+\sum_{j=3}^n\beta_j f_{n-j}.
\end{equation}
Put \(K'=\allowbreak -K\), \(z=\allowbreak K'\), \(g'=\allowbreak -g\); the reversed convolution uses
\(\bar f_n\). With \(G_m(z)=\allowbreak \sum_{j=\allowbreak 0}^{m-1}z^j\), \(G_0=\allowbreak 0\),
\(\|G_m\|\le\allowbreak 1\) for \(m=\allowbreak 0,\allowbreak 1\) or \(m\ge\allowbreak 3\), whereas
\(\|G_2\|<\allowbreak 6/5\). Frequencies have \(z=\allowbreak i\sqrt x\), \(0\le\allowbreak  x\le\allowbreak \theta\).
The full-mean inequality in Appendix~\ref{proof:Q5-C081} gives
\begin{equation}\tag{Q97.1.4}\label{eq:Q97:1.4}
\boxed{\E e^{8K'V}\le e^{8\E(K'V)+17\|V\|_2^2}}
\end{equation}
for every real \(V\), including \(V/C\) for any scalar \(C>\allowbreak 0\).
The bare-path bound of \eqref{eq:Q72:bare} is
\begin{equation}\tag{Q97.1.5}\label{eq:Q97:1.5}
\E[(p\bar f_n)^8]\le1-d/50\qquad(n\ge6).
\end{equation}
After zero-mass atoms are removed, \(p,\allowbreak q\ge\allowbreak \mu(x)>\allowbreak 0\).
The all-equal labelling gives positive tree messages. If \(d=\allowbreak 0\),
\(g=\allowbreak 0\), and leaf integration gives every message equal to one even
when \(K\ne0\). Divisions below use \(d>\allowbreak 0\).
Genuine equal pairs satisfy \(f_\ell+h_\ell=\allowbreak 2\E f_{\ell-1}\) and
\(0\le\allowbreak  f_\ell h_\ell\le\allowbreak 1\), and are restored by positive comparison.

\paragraph{Q97.2.1 The bare fourth path: cost at most \(3d\).}
Finite convolution gives the exact identity
\begin{equation}\tag{Q97.2.1}\label{eq:Q97:2.1}
\bar f_4=(1-d)q+zG_3(z)g'-d.
\end{equation}
Put \(C=\allowbreak 1-d\ge\allowbreak 2/3\), \(N=\allowbreak zG_3g'\). AM--GM on \(Cp\) and
\(\bar f_4\) gives
\fitdisplay{
p\bar f_4\le C\left(1+\frac{N-d}{2C}\right)^2.
}{}
The base is \((Cp+\bar f_4)/(2C)\ge\allowbreak 0\).
Using \(y^{16}\le\allowbreak  e^{16(y-1)}\) and \eqref{eq:Q97:1.4},
\begin{equation}\tag{Q97.2.2}\label{eq:Q97:2.2}
\log\E[(p\bar f_4)^8]\le8\log C+\frac8C(-d-\tau)
+\frac{17}{C^2}(d-\delta_2+\delta_3).
\end{equation}
Here \(\E N=\allowbreak -\tau\) and \(\|G_3g'\|^2=\allowbreak d-\delta_2+\delta_3\).
Let \(t=\allowbreak \delta_2/d\ge\allowbreak  d\). Since \(\delta_3\le\allowbreak \theta\delta_2\),
the squared norm divided by \(d\) is at most \(1-\kappa t\).
Use \(\log C\le\allowbreak -d\) and multiply the bound by \(C^2/d\):
\begin{equation}\tag{Q97.2.3}\label{eq:Q97:2.3}
-7+32d-8d^2-(67/32+8d)t\le-7+(957/32)d-16d^2.
\end{equation}
Subtracting \(3C^2\) from the last expression gives
\(-10+(1149/32)d-19d^2\). This increases on \([0,\allowbreak 1/3]\), since its
derivative is at least \(2231/96>\allowbreak 0\), and its endpoint value is
\(-41/288<\allowbreak 0\). Hence \eqref{eq:Q97:2.2} is at most \(3d\).

\paragraph{Q97.2.2 The bare fifth path: contraction.}
Retaining three convolution terms gives
\begin{equation}\tag{Q97.2.4}\label{eq:Q97:2.4}
\bar f_5=(1-2d)q+z(G_4(z)-d)g'-\tau.
\end{equation}
Put \(C=\allowbreak 1-2d\ge\allowbreak 1/3\), \(V=\allowbreak (G_4-d)g'\), \(N=\allowbreak zV\).
Then \(\E N=\allowbreak -\tau+d^2\), and the spectral blocks give
\fitdisplay{
\|V\|^2=(1-2d+d^2)d-(1-2d)\delta_2-\delta_3+\delta_4.
}{}
Using \(t=\allowbreak \delta_2/d\ge\allowbreak  d\), \(\delta_4\le\allowbreak \theta\delta_3\),
and \(\delta_3\ge\allowbreak \delta_2^2/d\),
\begin{equation}\tag{Q97.2.5}\label{eq:Q97:2.5}
\|V\|^2/d\le1-2d+d^2-Ct-\kappa t^2.
\end{equation}
The same weighted AM--GM and exponential inequality yield
\begin{equation}\tag{Q97.2.6}\label{eq:Q97:2.6}
\begin{aligned}
\log\E[(p\bar f_5)^8]\le{}&8\log C+\frac{8d}{C}(-2+2t+d)\\
&+\frac{17d}{C^2}[1-2d+d^2-Ct-\kappa t^2].
\end{aligned}
\end{equation}
The inequality \(\log C\le\allowbreak -2d-2d^2\), initially using only
\(-2d\), gives after multiplication by \(C^2/d\)
\(-15+70d-63d^2-Ct-17\kappa t^2\), decreasing for \(t\ge\allowbreak  d\).
Set \(t=\allowbreak d\) and include the remaining logarithmic term
\(-16dC^2\). The result is
\begin{equation}\tag{Q97.2.7}\label{eq:Q97:2.7}
P(d)=-15+53d-(227/32)d^2-64d^3.
\end{equation}
On \([0,\allowbreak 1/3]\), \(P'(d)\ge\allowbreak 431/16>\allowbreak 0\), and
\(P(1/3)=\allowbreak -425/864\). Thus the exponent is at most
\(-425d/(864C^2)\le\allowbreak -425d/864\), proving \eqref{eq:Q97:0.2}.

\paragraph{Q97.3 The triple envelope and the exact first step.}
Fix \(c\ge\allowbreak 5,\allowbreak L\ge\allowbreak 2\), \(F=\allowbreak \bar f_2f_3h_c\), and set
\begin{equation}\tag{Q97.3.1}\label{eq:Q97:3.1}
\begin{gathered}
B=\beta_{c-1},\quad m=c-3\ge2,\quad
W=\sum_{r=4}^cK^r\one=K^3G_m(K)g,\\
R=Bg-W,\qquad\Gamma=d+\beta_c+B.
\end{gathered}
\end{equation}
Since \(f_{c-1}\le\allowbreak  S_{c-1}-B\), positivity gives
\(h_c\le\allowbreak  TS_{c-1}-Bq\). Expanding and adding \(f_3\) yields
\begin{equation}\tag{Q97.3.2}\label{eq:Q97:3.2}
f_3+h_c\le2-\Gamma+R,\qquad0<F\le\bar f_2e^{R-\Gamma}.
\end{equation}
At \(c=\allowbreak 5\), \(B=\allowbreak d\) and \(\Gamma=\allowbreak 3d-\delta_2=\allowbreak 2d+\tau\ge\allowbreak 83d/32\).
For \(c\ge\allowbreak 6\), \(B\ge\allowbreak \tau\) and \(\Gamma\ge\allowbreak  d+2\tau\ge\allowbreak 35d/16\).
The exponential envelope follows by AM--GM on the two nonnegative arms.

First suppose \(0<\allowbreak d\le\allowbreak 1/8\). The row estimate and \(G_2\) bound give
\(|R|\le\allowbreak  d+(6/5)\theta\sqrt d<\allowbreak 1/8+(39/80)(3/8)<\allowbreak 1/2\).
The factorial bound \(n!\ge\allowbreak 2\cdot3^{n-2}\), \(n\ge\allowbreak 2\), bounds the
remainder by \(R^2/[2(1-|R|/3)]\le\allowbreak 3R^2/5\).
Write \(r=\allowbreak Kg\), \(G=\allowbreak g-r\), \(\bar f_2=\allowbreak 1-G\), so \(\E G^2=\allowbreak d+\delta_2\).
Keeping the complete mixed mean in the first row gives
\begin{equation}\tag{Q97.3.3}\label{eq:Q97:3.3}
\begin{gathered}
T(\bar f_2e^R)\le\bar f_3+zR+\delta,\\
\boxed{\delta=\E R-\E(gR)+\E(rR)
+\sqrt{(d+\delta_2)\E R^2}+\frac{12}{5}\E R^2.}
\end{gathered}
\end{equation}
Indeed, \(T(\bar f_2R)=\allowbreak \E R-\E(GR)-KR+K(GR)\),
\(K(GR)\le\allowbreak \E|GR|\le\allowbreak \sqrt{(d+\delta_2)\E R^2}\), and
\(T(\bar f_2R^2)\le\allowbreak 4\E R^2\).
For \(D_i=\allowbreak \E z^iR\), finite propagation is
\begin{equation}\tag{Q97.3.4}\label{eq:Q97:3.4}
e^\Gamma T^LF\le Y_L=\bar f_{L+2}+z^LR+\delta\bar f_{L-1}
+\sum_{i=1}^{L-1}D_i\bar f_{L-1-i}.
\end{equation}
The first perturbation is \(zR\), with coefficient one.

\paragraph{Q97.4.1 The first-step constant for \(c=\allowbreak 5\).}
Put \(t=\allowbreak \delta_2/d\), \(u=\allowbreak \delta_3/d\). The exact mixed moments are
\begin{equation}\tag{Q97.4.1}\label{eq:Q97:4.1}
\E R-\E(gR)+\E(rR)=-\delta_2-d^2+2\delta_3,\qquad
\E R^2=d^3+\delta_4+\delta_5-2d\delta_3.
\end{equation}
Since \(\theta(1+\theta)-2d>\allowbreak 0\) and \(u\le\allowbreak \theta t\),
\(\E R^2/d\le\allowbreak  J=\allowbreak d^2+[\theta^2(1+\theta)-2\theta d]t\).
Insert \(\sqrt{(1+t)J}\le\allowbreak (1+t)/6+(3/2)J\) into
\eqref{eq:Q97:3.3}:
\begin{equation}\tag{Q97.4.2}\label{eq:Q97:4.2}
\begin{aligned}
\delta/d\le{}&1/6-d+(39/10)d^2\\
&+t[-1/48+(39/10)\theta^2(1+\theta)-(39/5)\theta d].
\end{aligned}
\end{equation}
The \(t\)-coefficient is positive for \(d\le\allowbreak 1/8\). A lower bound,
using \(\theta^2>\allowbreak 4/25\), \(1+\theta>\allowbreak 7/5\), \(\theta<\allowbreak 5/12\), is
\(-1/48+(39/10)(4/25)(7/5)-(39/5)(5/12)(1/8)>\allowbreak 0\).
Set \(t\le\allowbreak \theta\). The resulting polynomial has derivative at most
\(-1-(39/5)\theta^2+39/40<\allowbreak 0\). Its value at zero is
\begin{equation}\tag{Q97.4.3}\label{eq:Q97:4.3}
81/512+(39/10)(98865/1048576)=5514615/10485760<17/32,
\end{equation}
with margin \(55945/10485760\). Thus \(\delta_+\le\allowbreak 17d/32\).

\paragraph{Q97.4.2 The first-step constant for \(c\ge\allowbreak 6\).}
Let \(e_j=\allowbreak \delta_j-\delta_{j+1}\). Then \(e_j\ge\allowbreak 0\) and
\(e_{j+1}\le\allowbreak \theta e_j\). Consecutive alternating sums give
\(\E W+\E(rW)\ge\allowbreak \kappa(\delta_2-\delta_3)\) and
\(\kappa\delta_3\le\allowbreak \E(gW)\le\allowbreak \delta_3\).
Together with \(B\ge\allowbreak \tau\ge\allowbreak \kappa d\),
\begin{equation}\tag{Q97.4.4}\label{eq:Q97:4.4}
\begin{gathered}
\E R-\E(gR)+\E(rR)\le-\kappa\delta_2+(1+\kappa)\delta_3-Bd,\\
\E R^2\le d^3+(\theta-2\kappa^2d)\delta_3.
\end{gathered}
\end{equation}
Here \(\theta-2\kappa^2d>\allowbreak 0\). Using \(u\le\allowbreak \theta t\), the mixed mean
divided by \(d\) is at most \((55/1024)t-\kappa d\), and
\(J=\allowbreak \E R^2/d\le\allowbreak  d^2+[\theta^2-2\kappa^2\theta d]t\).
Apply \(\sqrt{(1+t)J}\le\allowbreak (1+t)/8+2J\):
\begin{equation}\tag{Q97.4.5}\label{eq:Q97:4.5}
\begin{aligned}
\delta/d\le{}&1/8-\kappa d+(22/5)d^2\\
&+t[1/8+55/1024+(22/5)\theta^2-(44/5)\kappa^2\theta d].
\end{aligned}
\end{equation}
The \(t\)-coefficient is positive: dropping its first two positive
terms and using \(d\le\allowbreak 1/8\) leaves
\((22/5)\theta^2-(11/10)\kappa^2\theta
=(11/10)\theta(4\theta-\kappa^2)>0\).
After \(t\le\allowbreak \theta\), the derivative in \(d\) is at most
\(11/10-\kappa-(44/5)\kappa^2\theta^2<\allowbreak 0\).
The exact comparison is
\(81/160=\allowbreak 2654208/5242880<\allowbreak 2684396/5242880\).
The value at zero is \(80709/163840<\allowbreak 1/2\), with margin
\(1211/163840\). Hence \(\delta_+\le\allowbreak  d/2\) for every \(c\ge\allowbreak 6\).

\paragraph{Q97.4.3 All signed propagation tails.}
The positive means from \(-W\) have classes \(i\equiv2,\allowbreak 3\pmod4\),
with first terms at most \(\delta_3,\allowbreak \delta_4\). Those from \(Bg\)
have class \(i\equiv1\), first term \(Bd\). Increasing \(i\) by four
increases the moment index by two. Thus every finite positive sum obeys
\begin{equation}\tag{Q97.4.6}\label{eq:Q97:4.6}
\sum_{i\ge1}(D_i)_+\le\frac{\delta_3+\delta_4+Bd}{1-\theta^2}
\le\frac{169}{608}d+\frac{1024}{855}d^2.
\end{equation}
The individual bounds are
\begin{equation}\tag{Q97.4.7}\label{eq:Q97:4.7}
\begin{gathered}
-\delta_3\le D_1\le d^2,\quad0\le D_2\le\delta_3,\quad D_3\le\delta_4,\\
D_i\le\max(\delta_5,B\delta_3)\le\theta^4d<d/36\qquad(i\ge4).
\end{gathered}
\end{equation}
For the last bound, \(d\le\allowbreak 1/8<\allowbreak \theta^2\).
At \(i\equiv0\), both parts are nonpositive; class \(i\equiv1\)
starts at \(i=\allowbreak 5\) with \(B\delta_3\) and a nonpositive other part;
classes \(i\equiv2,\allowbreak 3\) start at \(i=\allowbreak 6,\allowbreak 7\) with
\(\delta_5,\allowbreak \delta_6\) and nonpositive \(Bg\) part.

\paragraph{Q97.5 Four retained path losses for \(L\ge\allowbreak 4\), \(d\le\allowbreak 1/8\).}
Here \(n=\allowbreak L+2\ge\allowbreak 6\). Retain the four nonnegative convolution terms
\(\tau\bar f_0,\allowbreak \tau\bar f_1,\allowbreak \tau\bar f_2,\allowbreak \tau\bar f_3\):
\begin{equation}\tag{Q97.5.1}\label{eq:Q97:5.1}
\bar f_n\le(1-3\tau)q
+z[G_{L+1}(z)-\tau(2+z)]g'-\tau(1-d).
\end{equation}
Indeed, \(\bar f_0+\bar f_1+\bar f_2+\bar f_3
=1-d+3q+2Kg-K^3\one\).
Set \(C=\allowbreak 1-3\tau>\allowbreak 0\), \(a=\allowbreak 1-2\tau\), and
\begin{equation}\tag{Q97.5.2}\label{eq:Q97:5.2}
P_{L,c}(z)=G_{L+1}(z)-\tau(2+z)-Bz^{L-1}-z^{L+2}G_m(-z).
\end{equation}
As \(R=\allowbreak -Bg'-z^3G_m(-z)g'\), the term \(z^LR\) in
\eqref{eq:Q97:3.4} combines with \eqref{eq:Q97:5.1} as
\(zP_{L,\allowbreak c}g'\). We prove
\begin{equation}\tag{Q97.5.3}\label{eq:Q97:5.3}
\boxed{\|P_{L,c}(K')\|\le1-2\tau=a.}
\end{equation}

\paragraph{Q97.5.1 The frequency certificate for \(L=\allowbreak 4\).}
At \(z=\allowbreak i\sqrt x\), \(0\le\allowbreak  x\le\allowbreak \theta\), write
\(G_m(-z)=\allowbreak U-zV\), where \(1-x\le\allowbreak  U,\allowbreak V\le\allowbreak 1\).
With \(v=\allowbreak 1-\tau\), the coordinates are
\fitdisplay{
\Re P=a-x+x^2+x^3U,\qquad
\Im P/\sqrt x=v-(1-B)x-x^3V.
}{}
They are nonnegative since \(a\ge\allowbreak 3/4\), \(v\ge\allowbreak 7/8\),
and \(x+x^3<\allowbreak 1/2\), and are at most \(a-x+x^2+x^3\) and
\(v-7x/8\), respectively. Squaring gives
\fitdisplay{
\begin{aligned}
|P|^2\le{}&a^2+(-1+2\tau+\tau^2)x+(5/4-9\tau/4)x^2\\
&+(49/64-4\tau)x^3-x^4+2x^5+x^6.
\end{aligned}
}{}
Since \(\tau\le\allowbreak 1/8\), the first coefficient is at most \(-47/64\),
and the next two positive coefficients are at most \(5/4,\allowbreak 49/64\).
Drop \(-x^4\) and divide by \(x\), treating \(x=\allowbreak 0\) by equality.
For \(x\le\allowbreak 5/12\), the total positive cost is at most
\fitdisplay{
\frac{25}{48}+\frac{1225}{9216}+\frac{625}{10368}
+\frac{3125}{248832}
=\frac{180800}{248832}<\frac{182736}{248832}=\frac{47}{64}.
}{}
Thus \(|P|^2\le\allowbreak  a^2\) at every frequency and for every \(c\).

\paragraph{Q97.5.2 The frequency certificate for \(L=\allowbreak 5\).}
The real coordinate is \(a-x+(1-B)x^2+x^4V\), and the imaginary
coordinate divided by \(\sqrt x\) is \(v-x+x^2+x^3U\), both nonnegative.
Bound these by \(a-x+x^2+x^4\) and \(v-x+x^2+x^3\), respectively:
\fitdisplay{
\begin{aligned}
|P|^2-a^2\le x[&-1+2\tau+\tau^2+(1-2\tau)x+(1-2\tau)x^2\\
&+(3-6\tau)x^3-3x^4+4x^5+x^6+x^7].
\end{aligned}
}{}
The derivative of the bracket in \(\tau\) is at least
\(2-2x-2x^2-6x^3\ge\allowbreak 37/96>\allowbreak 0\) for \(x\le\allowbreak 5/12\).
Take \(\tau=\allowbreak 1/8\) and drop \(-3x^4\). The remaining positive terms are
\((3/4)x+(3/4)x^2+(9/4)x^3+4x^5+x^6+x^7\).
Using \(x\le\allowbreak 5/12\), \(x^2<\allowbreak 7/40\), \(x^3<\allowbreak 3/40\),
\(x^5<\allowbreak 1/72\), \(x^6<\allowbreak 1/180\), \(x^7<\allowbreak 1/400\), their sum is below
\fitdisplay{
49/80+1/18+1/180+1/400=1217/1800<47/64.
}{}
This proves \(|P|^2\le\allowbreak  a^2\).

\paragraph{Q97.5.3 Every \(L\ge\allowbreak 6\).}
For comparison with the finite geometric remainder, define the scalar
rational function \(P_\infty(z)=\allowbreak 1/(1-z)-\tau(2+z)\). Directly,
\fitdisplay{
a^2-|P_\infty|^2=x\left[\frac{a}{1+x}-\tau^2\right]
\ge\frac{559}{1088}x.
}{}
The last coefficient follows from \(a\ge\allowbreak 3/4\), \(x\le\allowbreak 5/12\),
\(\tau\le\allowbreak 1/8\), giving \(9/17-1/64\).
Also \(a+|P_\infty|\le\allowbreak 2a\le\allowbreak 2\), so
\(|P_\infty|\le\allowbreak  a-559x/2176\le\allowbreak  a-x/4\).
The finite \(G_{L+1}\) remainder is at most
\(x^{(L+1)/2}/\sqrt{1+x}\le\allowbreak  x^{7/2}\).
The other two terms in \eqref{eq:Q97:5.2} are at most
\(x^{5/2}/8+(6/5)x^4\). After division by \(x\), their combined
bound, using \(x\le\allowbreak 5/12\), \(\sqrt x\le\allowbreak 2/3\), is
\fitdisplay{
25/216+5/144+25/288=205/864<1/4.
}{}
Hence \(|P_{L,\allowbreak c}|\le\allowbreak  a-11x/864\le\allowbreak  a\), with equality at \(x=\allowbreak 0\).
This completes \eqref{eq:Q97:5.3} for every finite \(L\ge\allowbreak 4,\allowbreak c\ge\allowbreak 5\).

\paragraph{Q97.5.4 Weighted AM--GM and the cost.}
Put \(P_{\rm cost}=\allowbreak 2\delta_++2\sum_{i\ge\allowbreak 1}(D_i)_+\), which bounds
the cost of every finite propagation. From \eqref{eq:Q97:3.4},
\eqref{eq:Q97:5.1} and \(\bar f_j\le\allowbreak 2\),
\fitdisplay{
0<e^\Gamma T^LF\le Cq+N+P_{\rm cost}-\tau(1-d),\qquad
N=zP_{L,c}(z)g'.
}{}
Here \(C=\allowbreak 1-3\tau\ge\allowbreak 5/8\),
\(\E N=\allowbreak -\beta_{n+1}+2\tau d+D_L\), and
\(\|P_{L,\allowbreak c}g'\|^2\le\allowbreak (1-2\tau)^2d\).
AM--GM on \(Cp\) and this nonnegative upper bound, followed by
\eqref{eq:Q97:1.4} applied to \(P_{L,\allowbreak c}g'/C\), gives
\begin{equation}\tag{Q97.5.4}\label{eq:Q97:5.4}
\begin{aligned}
\log\E[(pT^LF)^8]\le{}&-8\Gamma+8\log C\\
&+\frac8C[P_{\rm cost}-\tau(1-d)-\beta_{n+1}+2\tau d+D_L]\\
&+\frac{17}{C^2}(1-2\tau)^2d.
\end{aligned}
\end{equation}

\paragraph{Q97.5.5 The two-variable payment for \(c\ge\allowbreak 6\).}
The preceding bounds give
\(P_{\rm cost}\le\allowbreak (473/304)d+(2048/855)d^2\) and \(D_L\le\allowbreak  d/36\).
Put \(b=\allowbreak 1024/855\), \(A_0=\allowbreak 4333/2736\).
The bracket in \eqref{eq:Q97:5.4} is at most
\(A_0d+2bd^2-2\tau+3\tau d\), and \(\Gamma\ge\allowbreak  d+2\tau\).
Use \(\log C\le\allowbreak -3\tau\), multiply by \(C^2/d\), and set
\(s=\allowbreak \tau/d\in[\kappa,\allowbreak 1]\). Bound the positive costs by
\(A=\allowbreak 8/5\), \(b=\allowbreak 6/5\), obtaining
\begin{equation}\tag{Q97.5.5}\label{eq:Q97:5.5}
\begin{aligned}
F(s,d)={}&9+8A-56s+d[288s^2+(4-24A)s+16b]\\
&-d^2[360s^3+76s^2+48bs].
\end{aligned}
\end{equation}
For \(d\le\allowbreak 1/8\), the derivative of \(F_s\) in \(d\) is at least
\(538s-270s^2+4-24A-12b\ge\allowbreak 268s+4-24A-12b>\allowbreak 0\).
Its maximum in \(d\) is thus at \(1/8\); there it increases in
\(s\le\allowbreak 1\) and has value \(-11/4-3A-3b/4<\allowbreak 0\) at \(s=\allowbreak 1\).
Consequently \(F\) decreases in \(s\), with maximum at \(s=\allowbreak \kappa\).
At this \(s\), the linear \(d\)-coefficient exceeds \(75\), and the
absolute negative quadratic coefficient is below \(140\), using
\(\kappa^2>\allowbreak 1/3\), \(\kappa<\allowbreak 3/5\).
Hence the \(d\)-derivative exceeds \(75-140/4=\allowbreak 40\), and the maximum
is at \(d=\allowbreak 1/8\). Its exact value is
\begin{equation}\tag{Q97.5.6}\label{eq:Q97:5.6}
\begin{aligned}
F(\kappa,1/8)&=121/5-(306/5)\kappa+(557/16)\kappa^2-(45/8)\kappa^3\\
&=-\frac{1365979}{1310720}<-1.
\end{aligned}
\end{equation}
Thus \eqref{eq:Q97:5.4} is at most \(-d/C^2\le\allowbreak -d\).

\paragraph{Q97.5.6 The payment for \(c=\allowbreak 5\).}
Here \(P_{\rm cost}\) increases by \(d/16\), so
\(A_0=\allowbreak 563/342<\allowbreak 5/3\), and \(\Gamma=\allowbreak 2d+\tau\).
With \(A=\allowbreak 5/3\), \(b=\allowbreak 6/5\), the scaled expression is
\begin{equation}\tag{Q97.5.7}\label{eq:Q97:5.7}
\begin{aligned}
F_5(s,d)={}&1+8A-48s+d[240s^2+(52-24A)s+16b]\\
&-d^2[288s^3+148s^2+48bs].
\end{aligned}
\end{equation}
Its \(s\)-derivative increases in \(d\), because the corresponding
derivative is bounded below by
\(406s-216s^2+52-24A-12b\ge\allowbreak 190s+52-24A-12b>\allowbreak 0\).
At \(d=\allowbreak 1/8\), \((F_5)_s\) increases in \(s\le\allowbreak 1\), with value
\(3/8-3A-3b/4<\allowbreak 0\) at \(s=\allowbreak 1\). The maximum is again at \(s=\allowbreak \kappa\).
There the linear coefficient exceeds \(80\) and the absolute negative
quadratic coefficient is below \(152\); the \(d\)-derivative is thus
greater than \(80-152/4>\allowbreak 0\). At the endpoint,
\begin{equation}\tag{Q97.5.8}\label{eq:Q97:5.8}
\begin{aligned}
F_5(\kappa,1/8)&=251/15-(237/5)\kappa+(443/16)\kappa^2-(9/2)\kappa^3\\
&=-\frac{2547481}{983040}<-2.
\end{aligned}
\end{equation}
This row gives \(\E H^8\le\allowbreak  e^{-2d}\). Together the last two
calculations prove \(e^{-d}\) for all \(c\ge\allowbreak 5,\allowbreak L\ge\allowbreak 4,\allowbreak d\le\allowbreak 1/8\).

\paragraph{Q97.6 The two short connectors on the low domain.}
Fix \(\bar\delta=\allowbreak 17d/32\ge\allowbreak \delta\) and
\(\lambda=\allowbreak d-\bar\delta=\allowbreak 15d/32\). Replacing \(\delta\) by
\(\bar\delta\) increases the same positive upper bound. Put
\begin{equation}\tag{Q97.6.1}\label{eq:Q97:6.1}
Q_L(z)=G_{L+1}(z)-Bz^{L-1}-z^{L+2}G_m(-z).
\end{equation}

\paragraph{Q97.6.1 The connector \(L=\allowbreak 2\).}
The same \(U,\allowbreak V\in[1-x,\allowbreak 1]\) give
\(\Re Q_2=\allowbreak 1-x-x^2U\in[0,\allowbreak 1-x-x^2/2]\) and
\(\Im Q_2/\sqrt x=\allowbreak 1-B+x^2V\in[0,\allowbreak 1+x^2]\).
Thus \(|Q_2|^2\le\allowbreak 1-x+3x^3+x^4/4+x^5\le\allowbreak 1\), since after division
by \(x\), the positive cost is at most \(27/32\) for \(x\le\allowbreak 1/2\).
The exact convolution \(\bar f_4=\allowbreak S'_4-d(1+q)\) and finite propagation
give \(p+Y\le\allowbreak 2+c+zV_{\rm tot}\), with
\fitdisplay{
c=\bar\delta+D_1-2d,\qquad
V_{\rm tot}=Q_2(z)g'-\lambda\one,\qquad
\E(zV_{\rm tot})=-\tau+D_2.
}{}
By \eqref{eq:Q97:4.7}, \(d\le\allowbreak 1/8\) and \(\delta_3<\allowbreak d/6\),
\fitdisplay{
c+\E N\le[17/32+1/8+1/6-2-19/32]d=-85d/48\le-7d/4.
}{}
Since \(\E(Q_2g')=\allowbreak -d+D_1\),
\fitdisplay{
\|V_{\rm tot}\|^2\le d+\lambda^2+2\lambda(d-D_1)
\le d+(10/3)d^2\le17d/12<3d/2.
}{}
AM--GM and \eqref{eq:Q97:1.4} give exponent at most
\(-8(35/16)d-14d+(51/2)d=\allowbreak -6d\).
This proves the low-domain \(L=\allowbreak 2\) bound \(e^{-6d}\) for every \(c\ge\allowbreak 5\).

\paragraph{Q97.6.2 The connector \(L=\allowbreak 3\).}
The bound \(|G_4|=\allowbreak (1-x)\sqrt{1+x}\le\allowbreak 1-x/2-x^2/2\), together with
\(B\le\allowbreak 1/8\), \(\|G_m\|\le\allowbreak 6/5\) and \(\sqrt x\le\allowbreak 13/20\), gives
\(|Q_3|\le\allowbreak 1-3x/8+(7/25)x^2\le\allowbreak 1\).
Also \(\Re Q_3=\allowbreak 1-x+Bx-x^3V\ge\allowbreak 1-x-x^3>\allowbreak 1/2\).
Since \(\lambda\le\allowbreak 1/8\),
\(|Q_3-\lambda|^2\le\allowbreak |Q_3|^2-2\lambda\Re Q_3+\lambda^2\le\allowbreak 1\).
Use the exact convolution
\(\bar f_5=\allowbreak S'_5-\tau-dq-d\bar f_2\). Then
\fitdisplay{
\begin{gathered}
c=\bar\delta+D_1+D_2-2d-\tau,\quad
V_{\rm tot}=(Q_3-\lambda)g'-\eta\one,\\
\eta=2d-\bar\delta-D_1>0,\qquad
\E(zV_{\rm tot})=-\tau+\lambda d+D_3.
\end{gathered}
}{}
As \(D_1\le\allowbreak  d^2\), \(D_2\le\allowbreak \delta_3<\allowbreak d/6\), \(D_3\le\allowbreak \delta_4<\allowbreak d/14\),
\fitdisplay{
c+\E N\le[17/32+1/4+1/6+1/14-2-19/16]d
=-1457d/672\le-2d.
}{}
On the other hand, \(D_1\ge\allowbreak -\delta_3\) gives
\(\eta\le\allowbreak 157d/96<\allowbreak 5d/3\).
Also \(\E[(Q_3-\lambda)g']=\allowbreak -\tau+D_2\in[-\tau,\allowbreak 0]\), since
\(\tau\ge\allowbreak \kappa d>\allowbreak \delta_3\). Therefore
\fitdisplay{
\|V_{\rm tot}\|^2\le d+\eta^2+2\eta\tau
\le d+(55/9)d^2\le127d/72<9d/5.
}{}
The total exponent is at most
\(-35d/2-16d+153d/5=\allowbreak -29d/10<\allowbreak -2d\).
Thus every \(c\ge\allowbreak 5\) has the low-domain \(L=\allowbreak 3\) bound \(e^{-2d}\).

\paragraph{Q97.7 The high domain with the row mean retained.}
Let \(1/8\le\allowbreak  d\le\allowbreak 1/3\), keeping \(F=\allowbreak \bar f_2f_3h_c\). Set
\fitdisplay{
U=\sum_{r=3}^{c-1}K^r\one=K^2G_m(K)g,\qquad W=KU.
}{}
The finite alternating sum gives \(\E U=\allowbreak d-\beta_c\), so
\(B-\E U=\allowbreak \Gamma-2d\in[0,\allowbreak d]\).
The centered row estimate \eqref{eq:Q97:1.2} gives
\begin{equation}\tag{Q97.7.1}\label{eq:Q97:7.1}
\begin{gathered}
R=(\Gamma-2d)g-K(U-\E U),\\
|R|^2\le(\Gamma-2d)^2+\operatorname{Var}(U)\le d^2+\|U\|^2.
\end{gathered}
\end{equation}
The middle inequality is Cauchy on the two coordinates
\(g,\allowbreak \sqrt{1-g^2}\).
For \(c\ge\allowbreak 6\), \(m\ge\allowbreak 3\), and \(\|U\|^2\le\allowbreak \delta_3\le\allowbreak \theta^2d\).
Thus \(|R|/d\le\allowbreak \sqrt{1+8\theta^2}=\allowbreak \sqrt{297/128}<\allowbreak 25/16\).
Since \(\Gamma/d\ge\allowbreak 35/16\), \(\Gamma-R\ge\allowbreak 5d/8\).
For \(c=\allowbreak 5\), \(\|U\|^2=\allowbreak \delta_3+\delta_4\le\allowbreak \theta^2(1+\theta)d\),
so \(|R|/d\le\allowbreak \sqrt{11701/4096}<\allowbreak 7/4\).
Now \(\Gamma/d\ge\allowbreak 83/32\) yields
\(\Gamma-R\ge\allowbreak 27d/32\ge\allowbreak 5d/8\).
Equation \eqref{eq:Q97:3.2} therefore gives \(f_3h_c\le\allowbreak  e^{-5d/8}\)
for every \(c\ge\allowbreak 5\). Positivity implies
\(H\le\allowbreak  e^{-5d/8}p\bar f_{L+2}\).
Use the new \(n=\allowbreak 4\) cost for \(L=\allowbreak 2\), the new \(n=\allowbreak 5\) contraction for
\(L=\allowbreak 3\), and \eqref{eq:Q97:1.5} for \(L\ge\allowbreak 4\). Uniformly,
\begin{equation}\tag{Q97.7.2}\label{eq:Q97:7.2}
\E H^8\le e^{-2d}\qquad(L\ge2).
\end{equation}
At \(d=\allowbreak 1/8\) this meets the low-domain proof with the same direction.
The weakest low-domain bound is \(e^{-d}\), proving
\eqref{eq:Q97:0.1}. The zero boundary, genuine decorations and full
reversal are included by the finite positive comparison. Full reversal
replaces \(K\) by \(-K\) and reverses all arms, connector and receiving
unit arm together while preserving the measure and \(\delta_j\).

\paragraph{Q97.8 The ordinary catalogue for minimum arm at least two.}
For \(a<\allowbreak b_1<\allowbreak c_1<\allowbreak \cdots<\allowbreak b_k<\allowbreak c_k\), \(a\ge\allowbreak 2,\allowbreak k\ge\allowbreak 1\), the following
choices depend on the integer structure:
\fitdisplay{
\begin{array}{c|c|c|c}
\text{range}&F&Q(d)&\text{proof}\\\hline
b_1\ge5&\bar f_a\prod_i h_{b_i}f_{c_i}&e^{-d}&\ref{proof:Q5-C093}\\
b_1=4,\ a=2,3&\bar f_a\prod_i h_{b_i}f_{c_i}&e^{-d}
&k=1:\ \ref{proof:Q5-C095};\ k\ge2:\ \ref{proof:Q5-C093}\\
b_1=3,\ a=2,\ k\ge2&\bar f_2f_3h_{c_1}\prod_{i=2}^kh_{b_i}f_{c_i}
&e^{-d}&\ref{proof:Q5-C095}\\
(a,b_1,c_1)=(2,3,4),\ k=1&\bar f_2f_3h_4&e^{-3d/32}&\ref{proof:Q5-C095}\\
a=2,\ k=1,\ b_1=3,\ c_1\ge5,\ L=1&\bar f_2h_3f_{c_1}&e^{-d/27}&\ref{proof:Q5-C070}\\
a=2,\ k=1,\ b_1=3,\ c_1\ge5,\ L\ge2&\bar f_2f_3h_{c_1}&e^{-d}&\eqref{eq:Q97:0.1}
\end{array}
}{}
Except for the last two rows, the listed ranges have every \(L\ge\allowbreak 1\).
Strict integer order forces \(b_1=\allowbreak 3,\allowbreak 4\) or \(b_1\ge\allowbreak 5\), with \(a=\allowbreak 2\)
when \(b_1=\allowbreak 3\). The rows therefore give
\begin{equation}\tag{Q97.8.1}\label{eq:Q97:8.1}
\boxed{\E[\{E_rpT^L(E_tF)\}^8]\le e^{-d/27}\le1-d/28.}
\end{equation}
The last inequality follows from \(e^x\ge\allowbreak 1+x\) and \(27+d\le\allowbreak 28\).
Each row also retains its displayed stronger bound in its own direction.
